\begin{abstract}
Action Quality Assessment (AQA) busca cuantificar objetivamente la calidad de ejecución de acciones humanas en video, con aplicaciones en deporte, rehabilitación y educación. Los modelos de referencia se apoyan en arquitecturas 3D profundas (I3D, SlowFast) con alto coste computacional, lo que limita su uso en dispositivos con recursos limitados. Esta tesis aborda dos preguntas: (i) ¿en qué medida un \textit{pipeline} liviano basado en arquitecturas modernas (TSM-MobileNetV2, MobileNetV3), bien inicializado con pesos de ImageNet y entrenado con prácticas actuales (AdamW, cosine annealing, mixed precision, gradient accumulation), puede aproximarse al rendimiento de un Teacher 3D profundo?, y (ii) ¿qué componentes del \textit{pipeline} sostienen ese rendimiento, y dónde están sus límites de aplicabilidad?

Se propone un \textit{pipeline} reproducible aplicado a tres dominios de naturaleza distinta (AQA-7 multi-disciplina deportiva, MTL-AQA clavados especializados, JIGSAWS cirugía robótica), evaluado contra dos Teachers de referencia: I3D-R50 (referencia histórica) y SlowFast-R50 (referencia 3D moderna). En AQA-7, los Students livianos alcanzan SRCC de $0.9021 \pm 0.005$ (TSM-MobileNetV2) y $0.8907 \pm 0.005$ (MobileNetV3) sobre tres semillas, frente a $0.9052$ del I3D y $0.9158$ del SlowFast-R50, con brecha menor a $0.025$ SRCC y robusta entre semillas. Estos resultados utilizan sólo el $9\%$ de los FLOPs del Teacher I3D y reducen la latencia de inferencia a 40~ms por clip de 64 frames en una RTX 3060 Mobile. En MTL-AQA y JIGSAWS la brecha se mantiene menor a $0.01$ SRCC, confirmando estabilidad del \textit{pipeline} en dominios distintos.

Las ablaciones aíslan los componentes responsables del rendimiento: el preentrenamiento ImageNet aporta $\sim 0.031$ SRCC (sin él, el SRCC cae a $0.8713$) y el módulo Temporal Shift aporta $\sim 0.010$ SRCC adicional. Esto convierte el \textit{pipeline} en una propuesta interpretable, no en un efecto agregado opaco. La transferencia cruzada confirma límites: MTL-AQA $\to$ AQA-7 transfiere parcialmente (SRCC $\sim 0.55$); AQA-7 $\to$ JIGSAWS fracasa, delimitando el alcance del enfoque. Como análisis marginal de la destilación de conocimiento, se evalúa la variante \textit{+~KD} con pérdidas auxiliares de atención y alineación temporal; sólo 1 de 6 configuraciones (MobileNetV3 + AQA-7) mejora con destilación, mientras que las restantes degradan, en algunos casos de forma severa. Este resultado cuestiona la premisa heredada de la literatura de que la destilación es necesaria para cerrar la brecha entre Teacher y Student livianos. El código fuente, las configuraciones experimentales y los \textit{checkpoints} se publican como recurso reproducible.
\end{abstract}
